% arXiv submission. Compile: tectonic paper.tex
%   or: pdflatex paper && bibtex paper && pdflatex paper && pdflatex paper
% Form values are in ../ARXIV.md. The longer guide is whitepaper.tex.
\documentclass[11pt]{article}

\usepackage[T1]{fontenc}
\usepackage[utf8]{inputenc}
\usepackage[margin=1in]{geometry}
\usepackage{booktabs}
\usepackage{amsmath}
\usepackage{xurl}
\usepackage[round,authoryear]{natbib}
\usepackage[hidelinks]{hyperref}

\hypersetup{%
  pdftitle={NepalDTM: A National 8-Metre Bare-Earth Digital Terrain Model of Nepal},
  pdfauthor={Achyut Paudel}%
}

\title{NepalDTM: A National 8-Metre Bare-Earth Digital Terrain Model of Nepal\\
from Multi-Source Satellite Data Fusion}
\author{Achyut Paudel\\
\small \href{mailto:achyut.pdl@gmail.com}{achyut.pdl@gmail.com}}
\date{October 2026}

\begin{document}
\maketitle

\begin{abstract}
NepalDTM is an open bare-earth elevation model of Nepal on an 8~m grid.
Where the NASA High Mountain Asia 8~m stereo mosaic exists it is the surface; elsewhere the surface is Copernicus GLO-30.
A histogram gradient-boosting model, trained on 6,230,548 ICESat-2 and 7,035,904 GEDI ground points, predicts how far that surface sits above the ground.
On 4,463,463 held-out ICESat-2 points inside Nepal, leaving out whole 48~km blocks, the root-mean-square error is 6.10~m, against 9.07~m for Copernicus GLO-30, 8.47~m for GEDTM30, and 7.78~m for FABDEM.
The bias is $-0.15$~m.
Inside the delivered grid, 37\% of the area uses the 8~m mosaic and 63\% is filled from Copernicus, so the posted spacing is not 8~m terrain detail everywhere.
Tiles, the validation files, and the code that rebuilds the model are public.
\end{abstract}

\noindent\textbf{Keywords:} NepalDTM; digital terrain model; bare earth; ICESat-2; GEDI; Nepal

\section{Introduction}

Flood maps and landslide assessments need the elevation of the ground, not the elevation of trees and roofs.
Nepal has no public high-resolution bare-earth model.
Copernicus GLO-30 is a global 30~m radar surface \citep{esa2021copernicus}.
The High Mountain Asia 8~m mosaics are a stereo surface of glacier and snow regions \citep{shean2017hma}.
FABDEM \citep{hawker2022fabdem} and GEDTM30 \citep{ho2025gedtm30} remove vegetation and buildings from global 30~m data.
GEDTM30 is fit using ICESat-2, so a score against those points is not an independent test of that model.

NepalDTM keeps the 8~m mosaic where it exists, fills the rest from Copernicus, and subtracts a predicted above-ground height.
The correction is trained on ICESat-2 ATL08 \citep{markus2017icesat2,neuenschwander2019atl08} and GEDI L2A \citep{dubayah2020gedi}.
Heights are EGM2008 orthometric metres \citep{pavlis2012egm2008}.
The delivered grid covers 148,380~km\textsuperscript{2} inside the national outline used for tiling.

\section{Data and method}

The product is 324 tiles of 24~km on UTM zone 45N (EPSG:32645).
Each core is $3{,}000 \times 3{,}000$ pixels at 8~m.
A 1~km buffer is used while the predictors are built and is dropped from the delivered tile.

Table~\ref{tab:data} lists the inputs.
ICESat-2 points are kept when the height is valid, the segment is snow-free, the beam is strong, and the multiple-scattering warning is at most 1.
The usual mountain cut on height uncertainty is not used, because that threshold grows with slope.
GEDI points are the lowest-mode elevations that pass the SlideRule quality and degrade filters, come from a power beam, and have sensitivity above 0.95.
Each GEDI point has half the weight of an ICESat-2 point.
After those filters, 13,266,452 points fall on the tiles.

\begin{table}[ht]
  \centering
  \caption{Inputs. The HMA role is the base surface where the mosaic has a finite height. Copernicus fills the rest.}
  \label{tab:data}
  \small
  \begin{tabular}{@{}lllp{4.6cm}@{}}
    \toprule
    Dataset & Spacing & What it is & Role \\
    \midrule
    HMA DEM 8~m & 8~m & Stereo surface & Base surface \\
    Copernicus GLO-30 & 30~m & Radar surface & Fill, and a predictor \\
    ICESat-2 ATL08 & $\sim$20~m along track & Ground height & Training and test \\
    GEDI L2A & 25~m footprint & Ground under canopy & Training, half weight \\
    Canopy height & 1~m & Vegetation height & Predictor \\
    WorldCover 2021 & 10~m & Land-cover class & Predictor \\
    \bottomrule
  \end{tabular}
\end{table}

On each tile the mosaic is shifted from the ellipsoid onto EGM2008.
Pixels whose difference from Copernicus is more than 40~m from the median difference on that tile are dropped.
Gaps are filled from Copernicus.
The predictors at 8~m are the mean, 75th percentile, and maximum of the canopy-height map \citep{tolan2024canopy}; the majority WorldCover class \citep{zanaga2022worldcover}; slope, aspect, ruggedness, and topographic position at about 100~m and 500~m; the difference between the surface and Copernicus; a flag for mosaic versus Copernicus fill; and the surface height itself.

The target is surface minus ground.
A histogram gradient-boosting model \citep{pedregosa2011sklearn} with absolute-error loss predicts that residual.
Each fit uses at most 3,000,000 of the 13,266,452 points.
The bare-earth value is the surface minus the predicted residual.
A second model predicts the absolute error at each pixel.
The headline test is five-fold GroupKFold on 94 blocks of 48~km, so a test block is absent from its training fold.
A second test holds out whole satellite passes.
The single model that is written into the tiles is fit after those tests.

\section{Results}

Scores in Table~\ref{tab:national} are out-of-fold residuals on the 4,463,463 ICESat-2 points inside Nepal at which NepalDTM, Copernicus, GEDTM30, and FABDEM all have a value.
RMSE is the root-mean-square error.
NMAD is the normalised median absolute deviation, $1.4826$ times the median absolute deviation from the median.
Bias is the mean error; positive means the model is too high.
The same spatial predictions on all 6,230,548 ICESat-2 points, before the common-support cut, have RMSE 5.49~m.
That figure is not comparable with the other models.

\begin{table}[ht]
  \centering
  \caption{Held-out ICESat-2 error inside Nepal, in metres, on 4,463,463 points.}
  \label{tab:national}
  \begin{tabular}{@{}lrrr@{}}
    \toprule
    Model & RMSE & NMAD & Bias \\
    \midrule
    NepalDTM, 48~km blocks & 6.10 & 2.61 & $-0.15$ \\
    NepalDTM, satellite passes & 6.01 & 2.54 & $-0.13$ \\
    Copernicus GLO-30 & 9.07 & 3.51 & $+3.87$ \\
    GEDTM30 v1.2 & 8.47 & 3.96 & $+1.80$ \\
    FABDEM V1-2 & 7.78 & 3.14 & $+1.90$ \\
    \bottomrule
  \end{tabular}
\end{table}

Table~\ref{tab:zone} splits a larger common set of 6,227,829 points by the surface elevation.
Those rows include points in the tile buffer outside the national polygon.
Province scores, which do partition the 4,463,463 points inside Nepal, run from 1.78~m RMSE in Madhesh to 6.86~m in Karnali.
Under forest the common-set RMSE is 6.82~m, against 11.20~m for Copernicus.
Off forest it is 4.94~m, against 6.89~m.
Above 40 degrees of slope it is 10.81~m.
Where the base surface is the 8~m mosaic the common-set RMSE is 4.00~m; where it is Copernicus fill, 6.71~m.
Of the area inside the national outline, 37\% uses the mosaic and 63\% is Copernicus fill.

\begin{table}[ht]
  \centering
  \caption{Spatial-holdout RMSE in metres by surface elevation, on the 6,227,829-point common set.}
  \label{tab:zone}
  \begin{tabular}{@{}lrrrrr@{}}
    \toprule
    Zone & $n$ & NepalDTM & Copernicus & GEDTM30 & FABDEM \\
    \midrule
    Below 300~m & 1,251,581 & 1.96 & 4.53 & 2.12 & 2.46 \\
    300--1,000~m & 487,932 & 5.18 & 8.53 & 5.87 & 6.16 \\
    1--2.5~km & 1,076,095 & 6.24 & 8.69 & 7.57 & 7.33 \\
    2.5--4~km & 776,214 & 8.27 & 12.88 & 10.45 & 10.27 \\
    Above 4~km & 2,636,007 & 5.33 & 7.55 & 8.98 & 7.50 \\
    \bottomrule
  \end{tabular}
\end{table}

A separate check uses differential GPS surveyed on and beside Khumbu Glacier in 2015--2016 \citep{quincey2016khumbu}.
None of these points were used to train the model.
Against the \emph{surface} tile, the RMSE is 2.49~m at glacier margins ($n=83$), 2.96~m on the glacier surface ($n=388$), and 1.02~m on the wind-tower grid ($n=85$).
Against the bare-earth tile the same groups are 3.89~m, 6.01~m, and 6.59~m.
On the glacier surface the bare-earth model is 4.85~m low.
The correction is treating ice texture as something to remove.
The median absolute step across a delivered tile edge is 2.08~m, against 2.06~m inside a tile.
A shift search against ICESat-2 on open slopes finds no horizontal offset in any province.

\section{Limitations}

The 8~m spacing is the stored grid.
Terrain shape at that spacing exists where the HMA mosaic supplies the surface.
Elsewhere the hills and valleys are still the Copernicus 30~m shape, with a canopy adjustment on top.
Above 40 degrees the RMSE exceeds 10~m.
Part of that is the reference: a horizontal miss of a few metres becomes a large height error on a steep face.
Glacier values in the bare-earth model should be read as a corrected surface, not as bedrock.
The median predicted absolute error, stored on the error tile, is 4.3~m, and the 90th percentile is 8.8~m.

\section{Data and code}

Cloud Optimized GeoTIFF tiles (surface, bare earth, and predicted absolute error) are at
\url{https://pub-7a51429102ce456f85c0aee664338f8f.r2.dev/cog/}.
The web map reads
\url{https://pub-7a51429102ce456f85c0aee664338f8f.r2.dev/NepalDTM.pmtiles}.
The code that rebuilds the model, and the validation files this paper quotes, are at
\url{https://github.com/achyut/NepalDTM}.
From the \texttt{code} directory, \texttt{python run.py all} downloads the inputs, fits the correction, writes the tiles, and runs the validation.
A NASA Earthdata token in \texttt{\textasciitilde/.earthdata\_token} is required and is not in the repository.
The code writes three VRTs, \texttt{nepal\_dem\_8m.vrt}, \texttt{nepal\_dtm\_8m.vrt}, and \texttt{nepal\_aux\_8m.vrt}, whose sources are the local \texttt{out/} tiles.
Code is MIT.
The NepalDTM tiles are CC BY 4.0.
Input datasets keep their own licences.

\bibliographystyle{plainnat}
\bibliography{references}

\end{document}
